%% ---------------------------------------------------------------------
\section{La ley de Mason}
\label{sec:so-mason}
%% ---------------------------------------------------------------------

Con muchos lazos entrelazados, el álgebra de bloques se vuelve un juego de paciencia. Hay una alternativa que da la respuesta de una vez.

\begin{historia}[Samuel J.\ Mason y los grafos de flujo de señal]
En 1953, Samuel J.\ Mason, del MIT, propuso representar un sistema de ecuaciones lineales como un \emph{grafo de flujo de señal} y mostró cómo leer la retroalimentación en él~\citep{Mas53}; en 1956 publicó la fórmula general de la ganancia entre dos nodos, que tiene en cuenta los lazos que no se tocan~\citep{Mas56}. Su fórmula completó lo que Arnold Tustin había dejado sin resolver en un trabajo previo sobre diagramas de flujo. Hoy se la conoce como \emph{Mason's gain formula} o \clave{ley de Mason}.
\end{historia}

En un grafo de flujo, las señales son \emph{nodos} y cada \emph{rama} dirigida lleva una ganancia (una función de transferencia). Un nodo suma las señales de todas las ramas que le llegan y transmite esa suma por todas las que salen. Vocabulario~\citep{ogata10}:

\begin{definicion}[label=def:so-grafo]{Grafo de flujo de señal}
\begin{itemize}
  \item \textbf{Nodo de entrada} (\emph{source}): solo tiene ramas salientes. \textbf{Nodo de salida} (\emph{sink}): solo entrantes. \textbf{Nodo mixto}: ambas.
  \item \textbf{Trayecto}: un recorrido por ramas en el sentido de las flechas. \textbf{Trayecto directo}: va del nodo de entrada al de salida sin pasar dos veces por el mismo nodo.
  \item \textbf{Lazo}: un trayecto cerrado que no repite nodos. Dos lazos son \textbf{disyuntos} si no tienen ningún nodo en común (no se tocan).
  \item \textbf{Ganancia} de un trayecto o de un lazo: el producto de las ganancias de sus ramas.
\end{itemize}
\end{definicion}

\begin{definicion}[label=def:so-mason]{Ley de Mason}
La función de transferencia entre el nodo de entrada y uno de salida es
\begin{equation}
  T(s)=\frac{1}{\Delta}\sum_{k=1}^{n}P_k\,\Delta_k,
  \label{eq:so-mason}
\end{equation}
donde $n$ es el número de trayectos directos, $P_k$ la ganancia del $k$-ésimo, $\Delta$ el \clave{discriminante} del grafo,
\begin{equation}
  \Delta=1-\sum_a L_a+\sum_{b,c}L_bL_c-\sum_{d,e,f}L_dL_eL_f+\cdots
  \label{eq:so-discriminante}
\end{equation}
(uno, menos la suma de las ganancias de todos los lazos, más la suma de los productos de pares de lazos disyuntos, menos la de ternas de lazos disyuntos, y así sucesivamente), y $\Delta_k$ el discriminante de la parte del grafo que \emph{no toca} el trayecto $P_k$ (se calcula como $\Delta$, eliminando los lazos que tocan $P_k$).
\end{definicion}

\begin{cuidado}[Discriminante, no determinante]
En la versión anterior se llamaba ``determinante'' a $\Delta$. Lo llamamos \clave{discriminante} para no confundirlo con el determinante de una matriz (aunque, en el fondo, es el determinante del sistema de ecuaciones de los nodos, y por eso funciona).
\end{cuidado}

Antes de un caso grande, comprobemos con el lazo de siempre: un trayecto $P_1=G$, un lazo $L_1=-G\Hs$ que toca al trayecto, así que $\Delta=1+G\Hs$, $\Delta_1=1$ y $T=G/(1+G\Hs)$, la regla \eqref{eq:so-retro}.

\begin{ejemplo}[label=ej:so-mason]{Dos trayectos y cuatro lazos}
Halle $Y(s)/R(s)$ para el grafo de la \cref{fig:so-mason}.
\solucion \textbf{Trayectos directos:} $P_1=G_1G_2G_3G_4G_5G_6$ (nodos 1-2-3-4-5-6-7) y $P_2=G_1G_2G_7G_8G_9$ (nodos 1-2-3-8-9-7).

\textbf{Lazos:} $L_1=-G_2H_1$ (nodos 2-3), $L_2=-G_2G_3G_4H_4$ (2-3-4-5), $L_3=-G_4H_2$ (4-5) y $L_4=-G_8H_3$ (8-9).

\textbf{Lazos disyuntos:} en pares, $L_1L_3$, $L_1L_4$, $L_2L_4$ y $L_3L_4$ ($L_1$ y $L_2$ comparten los nodos 2 y 3; $L_2$ y $L_3$, el 4 y el 5). En ternas, solo $L_1L_3L_4$.
\[
  \Delta=1-(L_1+L_2+L_3+L_4)+(L_1L_3+L_1L_4+L_2L_4+L_3L_4)-L_1L_3L_4.
\]
\textbf{Cofactores:} $P_1$ toca los nodos 1 a 7, así que solo queda libre $L_4$: $\Delta_1=1-L_4$. $P_2$ toca 1, 2, 3, 8, 9 y 7, así que solo queda libre $L_3$: $\Delta_2=1-L_3$. Finalmente,
\[
  \frac{Y(s)}{R(s)}=\frac{P_1(1-L_4)+P_2(1-L_3)}{\Delta}.
\]
\end{ejemplo}

\begin{figure}[ht]
\centering
\includegraphics[width=0.95\linewidth]{cap03/mason_grafo}
\caption{Ejemplo~\ref{ej:so-mason}: grafo de flujo de señal con nueve nodos, dos trayectos directos y cuatro lazos (en coral, las ramas de retroalimentación).}
\label{fig:so-mason}
\end{figure}

\codigorepo{cap03_modelado_segundo_orden/C3_Ej6_Mason.m}{mason}{ejemplo \ref{ej:so-mason}: Mason contra resolver las ecuaciones de los nodos}

\begin{ideaclave}
El discriminante es siempre el denominador de la función de transferencia: \alerta{$\Delta(s)=0$ es la ecuación característica} y sus raíces son los polos del lazo cerrado. Para el lazo simple, $\Delta=1+G\Hs$: lo mismo que encontramos con el álgebra de bloques.
\end{ideaclave}

\begin{receta}[Aplicar la ley de Mason]
\begin{enumerate}
  \item Numerar los nodos. Listar \emph{todos} los trayectos directos con sus nodos.
  \item Listar todos los lazos con sus nodos y sus ganancias (incluyendo el signo de cada retroalimentación).
  \item Buscar pares, ternas\ldots\ de lazos que no compartan nodos y armar $\Delta$.
  \item Para cada trayecto, tachar los lazos que lo tocan y armar $\Delta_k$ con los que quedan.
  \item Aplicar \eqref{eq:so-mason} y, si se puede, verificar con álgebra de bloques o con \matlab.
\end{enumerate}
\end{receta}

\begin{pausa}
En la \cref{fig:so-mason}, si se elimina la rama $-H_4$, ¿cuánto vale el nuevo $\Delta$?
\respuesta{Quedan $L_1$, $L_3$ y $L_4$, todos disyuntos entre sí: $\Delta=1-(L_1+L_3+L_4)+(L_1L_3+L_1L_4+L_3L_4)-L_1L_3L_4=(1-L_1)(1-L_3)(1-L_4)$. Cuando ningún lazo se toca, el discriminante se factoriza.}
\end{pausa}
